\chapter{Introduction: a map of harmony}\label{ch:intro}

\section{Why a map?}
Musicians have always used spatial language for harmony: chords are ``close'' or ``distant'', a modulation
``travels'' to a ``remote'' key, a progression ``returns home''. The Tonnetz makes this language literal. It is a
diagram in which every point is a note and every small triangle is a chord, arranged so that geometric nearness
tracks a musical notion of nearness: two triangles that share an edge are two chords that share two notes and
can be connected by moving a single voice a short distance.

Once the picture is drawn, three branches of mathematics appear at once. As a \emph{graph} it invites questions
about paths and cycles: can we walk through all 24 major and minor chords, each exactly once, moving only between
neighbours? As a \emph{group action} it explains why transposing a progression preserves its shape. As a
\emph{surface} it has a topology: when the diagram is wrapped up so that each note appears only once, it becomes a
torus, and variants built from other chords become cylinders, spheres, M\"obius bands or even the real projective
plane. The goal of this report is to develop all three views carefully, to compute everything that can be
computed, and finally to test the map against real music.

\section{A short history}
\paragraph{Euler (1739).} Leonhard Euler's \emph{Tentamen novae theoriae musicae} \citep{euler1739} contains a
small lattice of note names in which horizontal neighbours differ by a perfect fifth and vertical neighbours by a
major third. For Euler, working in just intonation, the lattice was a way to organise frequency ratios built from
the primes 3 and 5; it was not yet a theory of chord progressions.

\paragraph{Nineteenth-century dualism.} Arthur von Oettingen \citep{oettingen1866} and Hugo Riemann
\citep{riemann1893} revived and extended the lattice (the German \emph{Tonnetz} is Riemann's usage). Their
harmonic ``dualism'' treated minor triads as mirror images of major ones, an idea that survives in the
inversional symmetry of Chapter~\ref{ch:prelims}.

\paragraph{Transformational and neo-Riemannian theory.} David Lewin's \emph{Generalized Musical Intervals and
Transformations} \citep{lewin1987} recast music theory in the language of groups acting on sets. Richard Cohn
then isolated three ``parsimonious'' transformations, $P$ (parallel), $L$ (leading-tone exchange) and $R$
(relative), that move a triad to another triad by shifting one voice by a semitone or a whole tone
\citep{cohn1996,cohn1997,cohn1998}. The Tonnetz became the natural stage for these operations, and a large body of
analysis of chromatic nineteenth-century music followed \citep{cohn2012}, extended later to popular music
\citep{capuzzo2004}.

\paragraph{Geometry and topology.} From the late 1990s mathematicians and theorists studied the Tonnetz as a
space. Douthett and Steinbach introduced the ``chicken-wire torus'', the dual graph on triads
\citep{douthett1998}; Crans, Fiore and Satyendra showed that the $PLR$ group is dihedral of order 24 and is the
dual of the transposition/inversion group \citep{crans2009}. Tymoczko connected chord spaces to orbifolds and
voice leading \citep{callender2008,tymoczko2011,tymoczko2012}, and Catanzaro classified the topology of every
``generalized Tonnetz'' built from a single trichord type in any equal temperament \citep{catanzaro2011}. Bigo and
collaborators represented chord collections as simplicial complexes and built the HexaChord software
\citep{bigo2013,bigo2016}. Yust \citep{yust2020}, Jevti\'c and \v{Z}ivaljevi\'c \citep{jevtic2021} and, most
recently, Boland and Hughston \citep{boland2025,boland2026} continue to generalise the construction, the latter
linking the Tonnetz to classical point--line configurations.

\paragraph{Computation and data.} In parallel, music information retrieval adopted the Tonnetz as a feature:
the tonal centroid of Harte, Sandler and Gasser \citep{harte2006} is implemented in the \code{librosa} library
\citep{mcfee2015} and used in chord recognition, and the Tonnetz has been used as an input representation for
neural networks \citep{chuan2018} and for genre classification from Tonnetz trajectories \citep{karystinaios2020}.
Large open chord corpora, notably ChoCo \citep{deberardinis2023} and Chordonomicon \citep{kantarelis2024}, now
make it possible to test Tonnetz-based hypotheses at scale; Xu, Hall and Rohrmeier's recent study of tonal
coherence on the Tonnetz across 2,800 pieces \citep{xu2026} is an example.

\section{What this report adds}
Beyond a unified tutorial exposition, the report contributes a set of verified, reproducible computations:
\begin{enumerate}
\item an independent enumeration of Hamiltonian cycles in the dual graph (62 undirected, 124 directed, as reported
by Albini and Antonini \citep{albini2009}) together with their classification into eight symmetry orbits and their
voice-leading costs (Chapter~\ref{ch:ham});
\item a machine-checked table of all twelve generalized Tonnetze in 12-tone equal temperament, including the
parity phenomenon that turns the chromatic strip into a M\"obius band for odd $N$ (Chapter~\ref{ch:topology});
\item a reconstruction of the presentation's whole-tone ``edge Tonnetz'' as a triangulated real projective plane,
contrasted with the actual pitch-class complex and with a Monte-Carlo census of edge-labelled gluings
(Chapter~\ref{ch:nonorient});
\item a four-corpus study with a shuffling null model (Chapter~\ref{ch:data}).
\end{enumerate}

\begin{keybox}[The one-sentence summary]
The Tonnetz is a triangulated torus whose vertices are the twelve notes and whose 24 triangles are the major
and minor triads; its dual is a 3-regular graph on those triads in which each edge is one of Cohn's
single-voice moves $P$, $L$ or $R$.
\end{keybox}

\begin{qa}
\Qu{In one sentence each, what are the three mathematical ``views'' of the Tonnetz used in this report?}
\A{As a graph (notes or chords joined by musical relations, studied through paths and cycles); as a group action
(transposition, inversion and the $P,L,R$ operations acting on chords); and as a surface (a triangulated torus
whose topology can be computed).}
\Qu{Why was Euler's original lattice not yet a torus?}
\A{Euler worked in just intonation, where no two lattice points have exactly the same frequency ratio, so the
lattice is an infinite plane. The torus appears only after identifying notes that are equal in 12-tone equal
temperament (Chapter~\ref{ch:tonnetz}).}
\Qu{Which three operations did Cohn call parsimonious, and what makes them parsimonious?}
\A{$P$, $L$ and $R$. Each keeps two notes of a triad fixed and moves the third by a semitone ($P$, $L$) or a
whole tone ($R$), so the voice leading is as small as possible between a major and a minor triad.}
\Qu{Name one audio-processing tool that uses the Tonnetz.}
\A{The tonal-centroid (``tonnetz'') feature of Harte, Sandler and Gasser, available as
\code{librosa.feature.tonnetz}; it projects a chroma vector onto three circles (fifths, minor thirds, major
thirds).}
\Qu{What is the practical difference between a tutorial review and a research article?}
\A{A tutorial review re-derives known results in a teachable order and points to the literature; this one also
adds reproducible code and a few new computations, but its main purpose is to make the field accessible.}
\end{qa}
